\documentclass{article}
\newcommand{\ip}[2]{(#1, #2)}
\begin{document}
\LaTeX\ is good at typesetting mathematical formulas
like
       \( x-3y + z = 7 \)
or
       \( a_{1} > x^{2n} + y^{2n} > x' \)
or
       \( \ip{A}{B} = \sum_{i} a_{i} b_{i} \).
The spaces you type in a formula are
ignored.  Remember that a letter like
       $x$
is a formula when it denotes a mathematical
symbol, and it should be typed as one.

Mathematical formulas may also be displayed.  A
displayed formula
is
one-line long; multiline
formulas require special formatting instructions.
   \[  \ip{\Gamma}{\psi'} = x'' + y^{2} + z_{i}^{n}\]
Don't start a paragraph with a displayed equation,
nor make one a paragraph by itself.
\end{document}
